
% =====================================================================================
% GUION Y FUENTES DEL CAPITULO (sep-2026)
%
% Este capitulo es el BLOQUE «Planteamiento de la comparativa» de la estructura oficial
% del tipo 3 (Instrucciones §2.6). La norma pide de el: identificar el problema concreto,
% las soluciones alternativas a evaluar, LOS CRITERIOS DE EXITO y LAS MEDIDAS QUE SE VAN A
% TOMAR. De ahi el orden: datos -> exploracion -> preprocesado -> modelos -> entrenamiento
% -> criterios y metricas.
%
% ⚠️ REPARTO CON EL CAPITULO 5 (decision del autor, sep-2026): TODO lo relativo a la
%    variable objetivo —Δ, el factor de supervivencia f, el umbral de señal y la eleccion
%    de los cinco observables— vive en el capitulo~\ref{cap:variable}, que va DELANTE.
%    Aqui se usa f dando por sabido lo que significa. No reabrir esa discusion.
%
% ⚠️ REPARTO CON EL CAPITULO 7: este capitulo DISEÑA (que se va a hacer y por que); el
%    capitulo~\ref{cap:comparativa} EJECUTA (que salio del barrido, hiperparametros
%    elegidos, curvas, incidencias). 🔴 NINGUNA CIFRA DE RENDIMIENTO ENTRA AQUI, ni de
%    validacion ni de test: si entra, el capitulo se contradice en cuanto se rehaga una
%    tabla, y los resultados quedan repartidos en dos sitios.
%
% FUENTES, por seccion:
%   6.1 -> TFM-Quantum/doc/dataset_info.md · doc/memoria/materiales.md §1 y §2
%   6.2 -> notebooks/01_eda/01a,01g,01h,01i · notebooks/01_eda_gnn/01j,01k
%   6.3 -> notebooks/02_preprocesado/DECISIONES_TOMADAS.md · src/preprocesado{,_gnn}.py
%   6.4 -> notebooks/04_gnn/NOTAS_ENTRENAMIENTO.md · doc/changes/2026-08_cambio_decisiones.md §9
%   6.5 -> notebooks/03_ridge_rf/NOTAS_ENTRENAMIENTO.md · notebooks/04_gnn/NOTAS_...md
%   6.6 -> notebooks/05_comparativa/METRICAS_Y_EVALUACION.md
%
% 🔴 CIFRAS PROHIBIDAS (son del dataset v1, obsoleto): «30-70 saltos» de distancia entre
%    nodos (la del v2 es 44), «42,7x» y «45,4x» del QFT (la citable es 3,69x), y
%    «0,344 -> 0,035» de la calidad de puerta al fijar el tamaño (en el v2 sube).
% =====================================================================================


\section{Creación del conjunto de datos}
\label{sec:dataset}
% =====================================================================================
% ⚠️ RECORTADO A TRES PAGINAS (sep-2026), por decision del autor. Lo que salio de aqui:
%
%   * el bloque «Lo que solo aparece en los datos reales» —las 17 puertas rotas los 42
%     dias, el 0,92 % de calibraciones que violan T2 <= 2·T1, la `rz` con error cero en
%     las 6.552 medidas y la correlacion 1,000000 entre `id`, `sx` y `x`—. Es anatomia
%     del chip, no planteamiento de la comparativa. Si se repone, su sitio es un anexo.
%   * el detalle de los 1.952 registros de calibracion por dia frente a los 976 del
%     conjunto base.
%   * la decision sobre las puertas averiadas, que baja de parrafo a frase.
%
% Las cifras de todo lo retirado siguen en doc/auditoria_dataset.md y en
% notebooks/00_anatomia_del_dataset/.
% =====================================================================================

Para la obtención del conjunto de datos, la única vía posible fue fabricarlo mediante simuladores. Las razones son las siguientes:

\begin{itemize}
    \item \textbf{Cuota.} El plan abierto de IBM concede diez minutos de procesador al mes
          \cite{ibm2026productos}.
          Etiquetar veinte mil circuitos con ese presupuesto no era posible de forma gratuita.
    \item \textbf{Latencia.} Los trabajos entran en una cola compartida, y veinte mil
          ejecuciones con su espera no caben en el calendario establecido.
    \item \textbf{La etiqueta no existe en el hardware.} El modelo tiene que aprender la \emph{diferencia} entre el resultado ideal y el ruidoso, y una máquina cuántica \textbf{solo devuelve el segundo} (Si el valor exacto se pudiera medir en la máquina, no habría nada que mitigar).
\end{itemize}

De ahí el planteamiento: un \textbf{gemelo digital}. Cada circuito se simula dos veces, una sin
ruido y otra con el ruido del procesador reconstruido a partir de su calibración real de ese
día. El valor ideal no se estima, se calcula; y el ruido no se inventa, se reconstruye.

Este planteamiento fija el tamaño máximo de los circuitos. Simular el estado de $n$ qubits exige seguir $2^n$ amplitudes: unos $4{,}3$~GB para catorce qubits y $17{,}2$~GB para quince, \textbf{por proceso}. El tope que se establece es de \textbf{quince qubits} debido a la memoria disponible.



El procesador imitado es \texttt{ibm\_kingston}, de la familia \textbf{Heron r2}: $156$ qubits
y una puerta de dos qubits nativa que es la \textbf{CZ}, simétrica, sobre un conjunto base
$\{$\texttt{cz}, \texttt{id}, \texttt{rz}, \texttt{sx}, \texttt{x}$\}$. Todos los circuitos se
transpilan a ese conjunto, así que lo que el modelo ve no son las puertas que escribió el
usuario, sino las que la máquina ejecutaría de verdad.

De ese procesador se descargaron \textbf{$42$ días completos} de calibración, del 13 de junio
al 27 de julio de 2026. Cada día aporta, por qubit, su $T_1$, su $T_2$, el error de lectura en
las dos direcciones y la duración de la medida; y por cada puerta física del chip, su error y
su duración reales.

% ⚠️ FORMATO: titulo arriba y fuente abajo (Instrucciones §1.3, p. 6).
% Figura propia, de notebooks/00_anatomia_del_dataset/, copiada desde
% TFM-Quantum/figures/00_anatomia/chip/04_drift_42_dias.pdf.
\begin{figure}[htbp]
    \centering
    \caption{Recalibración diaria del procesador a lo largo de los $42$ días.}
    \label{fig:drift_t1}

    \includegraphics[width=0.92\textwidth]{cap06/drift_t1}

    {\footnotesize Fuente: elaboración propia.}
\end{figure}

Falta la colocación (véase el panel izquierdo de la Figura~\ref{fig:chip}). Cada circuito se sitúa sobre una \textbf{región conexa aleatoria} de qubits físicos, distinta en cada muestra, en lugar de usar siempre los mismos. Así el modelo ve una variedad amplia de calidades de qubits. Del muestreo se excluye un solo qubit, el $146$, que no tiene $T_1$ ni $T_2$ medidos ninguno de los
$42$ días: está averiado, y su telemetría no es un dato sino un hueco.


Las puertas que el chip reporta como averiadas \textbf{se mantienen} ya que es un fenómeno usual y aprendible.
% ⚠️ FORMATO: titulo arriba y fuente abajo (Instrucciones §1.3, p. 6).
% Dos figuras propias, de notebooks/00_anatomia_del_dataset/00b, copiadas desde
%   figures/00_anatomia/chip/12_mapa_averias_memoria.pdf
%   figures/00_anatomia/chip/13_region_memoria.pdf
% Las dos SIN titulo dentro de la imagen: lo pone el \caption.
% ⚠️ A media anchura los numeros de qubit NO se leen, y es una decision consciente del
%    autor (sep-2026): la figura vale para ver DONDE estan las averias y DONDE cae la
%    region, no para localizar un qubit concreto. No «arreglarlo» agrandandolas.
\begin{figure}[htbp]
    \centering
    \caption{El chip por dentro. A la izquierda, la región
             conexa sobre la que se coloca un circuito concreto. A la derecha, las averías y la invalidez del qubit $146$.}
    \label{fig:chip}

    \begin{minipage}[t]{0.49\textwidth}
        \centering
        \includegraphics[width=\textwidth]{cap06/region_chip}
    \end{minipage}
    \begin{minipage}[t]{0.49\textwidth}
        \centering
        \includegraphics[width=\textwidth]{cap06/mapa_averias}
    \end{minipage}\hfill
    

    {\footnotesize Fuente: elaboración propia.}
\end{figure}


El conjunto tiene \textbf{$20{.}000$ muestras} y $18{.}147{.}608$ nodos, repartidas en
$16{.}000$ de entrenamiento y validación y $4{.}000$ de prueba. Cada muestra es un circuito de
entre $5$ y $15$ qubits y entre $5$ y $50$ capas lógicas de profundidad. Los dos extremos de la
profundidad están elegidos por la misma razón: por debajo de cinco capas apenas hay ruido que
medir, y por encima de cincuenta la decoherencia ha destruido ya toda la señal.

% ⚠️ FORMATO: titulo arriba y fuente abajo (Instrucciones §1.3, p. 6).

El reparto entre entrenamiento y prueba se organiza para que el conjunto de prueba obligue a generalizar en \textbf{tres
direcciones independientes} (Figura~\ref{fig:tres_ejes}):

% FIGURA PROPIA en TikZ. Codigo en figuras/cap06/tres_ejes_tikz.tex.
% ⚠️ SUSTITUYE a dos cosas: a la antigua tabla `tab:ejes`, que decia lo mismo en prosa
%    numerada, y al PNG `02_los_tres_ejes` del cuaderno 00b, que eran tres cajas en fila.
%    El abanico dice ademas que los tres ejes SALEN del mismo conjunto.
\begin{figure}[htbp]
    \centering
    \caption{Cómo se reparte el conjunto, y de dónde salen los tres ejes de generalización.}
    \label{fig:tres_ejes}

    % =====================================================================================
% Cómo se reparte el conjunto, y de dónde salen los tres ejes — figura propia en TikZ.
%
% CUATRO NIVELES:
%   1. el conjunto entero;
%   2. entrenamiento, validación y prueba, con QUÉ HAY dentro de los dos primeros;
%   3. de PRUEBA salen los tres ejes de generalización, porque es ahí —y solo ahí— donde
%      aparecen tipos, tamaños y días que el modelo no ha visto;
%   4. de TIPO y de TAMAÑO sale además cómo se reparte por dentro cada uno de esos ejes.
%      TIEMPO no abre cuarto nivel: sus doce días no se subdividen.
%
% ⚠️ SIN TITULO dentro de la imagen: lo pone el \caption de LaTeX (norma del TFE, §1.3).
%
% ⚠️ ANCHOS. Los cuatro niveles tienen que caber en los 15 cm de la caja de texto, así que
%    cada columna está medida: 1,9 + 3,7 + 2,3 + 2,9 cm de texto más sus márgenes suman
%    unos 14,7. NO ensanchar ninguna sin estrechar otra, o la figura se sale de la caja.
%    Los niveles 3 y 4 van en \scriptsize por el mismo motivo.
%
% COLORES. Los tres ejes son categorías hermanas sin orden ni significado físico, así que
% aquí NO aplica el código de color de la memoria (azulunir = el objeto, acentofig = lo
% que se le hace). Se toman los tres tonos de la paleta de tipos de circuito del proyecto
% —src/viz/estilo.py::CIRCUITOS—, ya validados contra daltonismo y contraste. El cuarto
% nivel hereda el color de su eje, para que se vea de quién cuelga.
%
% ⚠️ \shorthandoff{<>} es obligatorio: `babel` en español hace activos < y >, y eso rompe
%    la sintaxis de puntas de flecha de TikZ.
% =====================================================================================
\begingroup
\shorthandoff{<>}%
% Dentro de cajas tan estrechas, LaTeX parte «qubits» y «posteriores». Se prohíbe.
\hyphenpenalty=10000\exhyphenpenalty=10000\relax

\definecolor{ejetipo}{HTML}{E34948}    % rojo    · el eje de tipo de circuito
\definecolor{ejetam}{HTML}{1BAF7A}     % verde   · el eje de tamaño
\definecolor{ejetiempo}{HTML}{4A3AA7}  % índigo  · el eje temporal

\begin{tikzpicture}[
    >={Stealth[length=2mm]},
    font=\footnotesize,
    origen/.style = {draw=azulunir, line width=1pt, rounded corners=3pt,
                     fill=azulunir!6, align=center, inner sep=5pt, text width=1.9cm},
    grupo/.style  = {draw=#1, line width=0.9pt, rounded corners=3pt, fill=#1!5,
                     align=left, inner sep=5pt, text width=3.7cm},
    eje/.style    = {draw=#1, line width=0.9pt, rounded corners=3pt, fill=#1!6,
                     align=left, inner sep=4pt, text width=2.3cm, font=\scriptsize},
    hoja/.style   = {draw=#1!55, line width=0.7pt, rounded corners=2pt, fill=#1!3,
                     align=left, inner sep=4pt, text width=2.9cm, font=\scriptsize},
    flecha/.style = {->, line width=1pt, draw=#1},
    ramita/.style = {->, line width=0.7pt, draw=#1!65},
  ]

  % --- nivel 1: el conjunto -----------------------------------------------------------
  \node[origen] (datos) at (0,0.6)
    {\textbf{El conjunto\\de datos}\\[2pt]
     {\scriptsize\color{gray} $20{.}000$}};

  % --- nivel 2: los tres grupos -------------------------------------------------------
  \node[grupo=azulunir]  (train) at (3.6, 4.2)
    {\textbf{Entrenamiento}\\[2pt]
     $13{.}344$ circuitos\\
     días $0$--$24$\\
     aleatorios · $5$--$11$ qubits};

  \node[grupo=azulunir]  (val)   at (3.6, 1.7)
    {\textbf{Validación}\\[2pt]
     $2{.}656$ circuitos\\
     días $25$--$29$\\
     aleatorios · $5$--$11$ qubits};

  \node[grupo=acentofig] (test)  at (3.6, -1.4)
    {\textbf{Prueba}\\[2pt]
     $4{.}000$ circuitos\\
     días $30$--$41$};

  \draw[flecha=azulunir]  (datos.east) to[out=45,  in=180] (train.west);
  \draw[flecha=azulunir]  (datos.east) to[out=15,  in=180] (val.west);
  \draw[flecha=acentofig] (datos.east) to[out=-45, in=180] (test.west);

  % --- nivel 3: los tres ejes, colgando de PRUEBA --------------------------------------
  \node[eje=ejetipo]   (tipo)   at (8.0,  1.8) {\textbf{\color{ejetipo}TIPO}};
  \node[eje=ejetam]    (tam)    at (8.0, -1.4) {\textbf{\color{ejetam}TAMAÑO}};
  \node[eje=ejetiempo] (tiempo) at (8.0, -4.0)
    {\textbf{\color{ejetiempo}TIEMPO}\\doce días posteriores};

  \draw[flecha=ejetipo]   (test.east) to[out=35,  in=180] (tipo.west);
  \draw[flecha=ejetam]    (test.east) to[out=0,   in=180] (tam.west);
  \draw[flecha=ejetiempo] (test.east) to[out=-40, in=180] (tiempo.west);

  % --- nivel 4: cómo se reparte por dentro cada eje ------------------------------------
  % TIEMPO no abre nivel: los doce días no se subdividen.
  \node[hoja=ejetipo] (tipoA) at (11.9,  2.7) {$30$\,\% aleatorios};
  \node[hoja=ejetipo] (tipoB) at (11.9,  0.9) {$70$\,\% seis circuitos específicos nuevos};

  \node[hoja=ejetam]  (tamA)  at (11.9, -0.6) {de $5$ a $11$ qubits};
  \node[hoja=ejetam]  (tamB)  at (11.9, -2.3) {de $12$ a $15$ qubits};

  \draw[ramita=ejetipo] (tipo.east) to[out=30,  in=180] (tipoA.west);
  \draw[ramita=ejetipo] (tipo.east) to[out=-30, in=180] (tipoB.west);
  \draw[ramita=ejetam]  (tam.east)  to[out=30,  in=180] (tamA.west);
  \draw[ramita=ejetam]  (tam.east)  to[out=-30, in=180] (tamB.west);

\end{tikzpicture}%
\endgroup


    {\footnotesize Fuente: Claude Code \cite{anthropic_claude}.}
\end{figure}

\begin{itemize}
    \item \textbf{Tipo.} El entrenamiento es aleatorio al $100$\,\% (del mismo modo validación) y en prueba se incluyen, a parte de circutios aleatorios, seis circuitos con estructura diferente escogida por elección frecuente en la literatura (véase la Tabla~\ref{tab:tipos}). Si el modelo acierta en ellos, es porque aprendió física del ruido y no patrones de circuito.
    \item \textbf{Tamaño.} El entrenamiento llega hasta once qubits y la prueba hasta quince.
    \item \textbf{Tiempo.} Los días de calibración de los dos grupos son \textbf{disjuntos}: del 13 de junio al 14 de julio para entrenar, del 15 al 27 de julio para probar. El modelo se evalúa sobre un estado del chip que no vio al entrenar.
\end{itemize}

El conjunto de prueba consta de un $30$\,\% de
circuitos aleatorios que se divide por la mitad
entre circuitos de hasta once qubits y circuitos de doce a quince. Los primeros permiten medir la deriva temporal con el tipo y el tamaño fijos; los segundos, la extrapolación en
tamaño con el tipo fijo. El resto del conjunto de prueba son 6 circuitos específicos, elegidos por frecuencia en la literatura, y divididos de la misma forma por número de qubits. En la Tabla~\ref{tab:tipos} aparece la justificación de la elección de estos cirucitos por la literatura.

\begin{table}[htbp]
    \centering
    \caption{Los siete tipos de circuito y la razón de que estén.}
    \label{tab:tipos}
    {\footnotesize
    Fuente: elaboración propia. \begin{tabular}{@{}llp{0.52\textwidth}@{}}
        \toprule
        \textbf{Tipo} & \textbf{Dónde aparece} & \textbf{Por qué está} \\
        \midrule
        Aleatorio & entrenamiento y prueba
            & Es la física del ruido sin la estructura de ningún algoritmo. Es el $100$\,\% del
              entrenamiento. \\
        HEA & solo prueba
            & \textit{Hardware-Efficient Ansatz}, el circuito de referencia de los algoritmos
              variacionales. \\
        TFIM & solo prueba
            & Simulación de un hamiltoniano de Ising con campo transverso; referencia estándar
              en los bancos de prueba de mitigación. \\
        QAOA & solo prueba
            & El algoritmo de optimización combinatoria más citado del estado del arte. \\
        QFT & solo prueba
            & Transformada de Fourier cuántica. \textbf{Aportación propia}: ningún trabajo de
              mitigación por aprendizaje automático la evalúa de forma sistemática. \\
        GHZ & solo prueba
            & El circuito de entrelazamiento por excelencia, y un caso extremo: su
              magnetización media es nula por construcción. \\
        Bernstein--Vazirani & solo prueba
            & Salida determinista, un régimen de ruido completamente distinto al del circuito
              aleatorio. \\
        \bottomrule
    \end{tabular}}
\end{table}

% =====================================================================================
\section{Análisis exploratorio de los datos}
\label{sec:eda}
% ⚠️ SECCION BREVE A PROPOSITO (sep-2026, decision del autor): aqui solo se enuncia QUE se
%    buscaba en el analisis. Los hallazgos concretos se citan mas adelante, cada uno en el
%    punto donde decide algo, en lugar de listarlos dos veces. El desarrollo completo esta
%    en notebooks/01_eda/ y en notebooks/02_preprocesado/DECISIONES_TOMADAS.md.

Antes de fijar un preprocesado o una arquitectura se hizo un análisis descriptivo del conjunto de datos.
Sus objetivos fueron:

\begin{itemize}
    \item \textbf{Localizar las casillas sin etiqueta} y decidir qué hacer con ellas.
    \item \textbf{Análsis de relación entre las variables y el factor de supervivencia.}
    \item \textbf{Reducir la dimensionalidad} descartando las columnas redundantes.
    \item \textbf{Comprobación de escalas y asimetrías}, que son las que deciden qué
          transformaciones hacen falta antes de entrenar.
    \item \textbf{Análisis de la información que sobrevive dentro de un mismo circuito}, porque es
          la única que el grafo puede aprovechar y el resumen tabular no.
    \item \textbf{Análisis descriptivo de la forma de los grafos} 
\end{itemize}

De ese análisis salen las decisiones de preprocesado del apartado~\ref{sec:prep} y del
diseño de los modelos del capítulo~\ref{cap:comparativa}. Los apartados que siguen recogen lo que
se concluyó en cada uno de esos frentes.

\subsection{Observables sin etiqueta}
\label{sec:eda_nulos}

El umbral de señal del capítulo~\ref{cap:variable} deja sin etiqueta entre el $4{,}2$\,\% y el
$14{,}5$\,\% de las casillas, según el observable, y la paridad es con diferencia la más
castigada. Aun así, el $69{,}6$\,\% de los circuitos de entrenamiento conserva los cinco
objetivos completos y ninguno se queda sin ninguna etiqueta. 

La idea de imputar las varibles se descartó porque los valores que eliminaba el umbral, fijado en $0.002$, eran datos que no se podían deducir de ningún otro porque a partir de ese umbral no existia ningún dato que pudiera deducir.

\textbf{Decisión: no se imputa
nada}, y se entrena de la siguiente forma:
\begin{itemize}
    \item Cinco modelos Ridge, uno por cada observable.
    \item Cinco modelos Random Forest, uno por cada observable.
    \item Tres modelos distintos multi-salida GNN con máscara booleana.
\end{itemize}

\subsection{Relación de cada familia de variables con la degradación}
\label{sec:relaciones}

La familia que más predice el tiempo del circuito medido frente a la decoherencia, con tres columnas.

De todas las familias, la única que no aporta apenas valor es  la telemetría del hardware con trece columnas. 

\textbf{Decisión: se quedan todas}, porque no afecta ni a Ridge ni a Random Forest, y porque
dentro del grafo sí varían (apartado~\ref{sec:variacion_intra}).

% Figura propia, de notebooks/01_eda/01h, copiada desde figures/eda/relaciones/02_familias.pdf.
\begin{figure}[htbp]
    \centering
    \caption{Cuánto predice cada familia de variables, por observable, medido de dos
             formas.}
    \label{fig:familias}

    \includegraphics[width=\textwidth]{cap06/familias}

    {\footnotesize Fuente: elaboración propia.}
\end{figure}

\subsection{Telemetría de un mismo circuito y forma de los grafos}
\label{sec:variacion_intra}

La correlación del apartado anterior se mide sobre la fila, que promedia el circuito
entero. Falta por saber qué sobrevive a ese promedio, ya que es lo que separa a los
dos modelos: el resumen tabular solo ve la media y el grafo ve nodo a nodo.

\textbf{La telemetría del qubit sigue variando dentro de un mismo circuito}
(Figura~\ref{fig:variacion_intra}) y es la única señal que el grafo puede aprovechar y la tabla no.

% ⚠️ FORMATO: titulo arriba y fuente abajo (Instrucciones §1.3, p. 6).
% Figura propia, de notebooks/01_eda_gnn/01j, copiada desde
% figures/eda_gnn/nodos/05_variacion_intra_memoria.pdf.
% ⚠️ Es la version de memoria: sin titulo, sin las dims que no varian, con nombres legibles
%    y con la telemetria RESALTADA, que es el argumento del apartado.
\begin{figure}[htbp]
    \centering
    \caption{Cuánto varía la telemetría del qubit dentro de un mismo circuito.}
    \label{fig:variacion_intra}

    \includegraphics[width=\textwidth]{cap06/variacion_intra}

    {\footnotesize Fuente: elaboración propia.}
\end{figure}

Los grafos de este conjunto son \textbf{muy profundos}
(Figura~\ref{fig:forma_grafos}): mediana de $805$ nodos, mínimo de $18$ y máximo de $5{.}056$, con las aristas creciendo de forma
lineal con los nodos. Ese factor decide una cosa del modelo de grafos: cualquier resumen que
\emph{sume} los nodos mediría el tamaño del circuito en lugar de su contenido
(apartado~\ref{sec:gnn}).

% Figura propia, de notebooks/01_eda_gnn/01k, copiada desde
% figures/eda_gnn/grafos/01_tamano.pdf.
\begin{figure}[htbp]
    \centering
    \caption{Tamaño de los grafos del conjunto de entrenamiento. La línea marca la mediana.}
    \label{fig:forma_grafos}

    \includegraphics[width=\textwidth]{cap06/forma_grafos}

    {\footnotesize Fuente: elaboración propia.}
\end{figure}

\subsection{Colinealidad}
\label{sec:colinealidad}

\begin{figure}[htbp]
    \centering
    \caption{Correlación entre las $54$ variables candidatas, antes de descartar nada.}
    \label{fig:colinealidad}

    \includegraphics[width=\textwidth]{cap06/colinealidad}

    {\footnotesize Fuente: elaboración propia.}
\end{figure}


De las $54$ columnas candidatas
(Figura~\ref{fig:colinealidad}) hay $11$ \textbf{identidades exactas} y numerosos\textbf{bloques redundantes} con
correlaciones de hasta $0{,}999$.

\textbf{Decisión: se descartan $20$ columnas y quedan $34$}, conservando un representante de
cada bloque. La matriz que queda después está en el apartado~\ref{sec:sol_ridge}, junto al
primer modelo que la consume.

% Figura propia, de notebooks/01_eda/01g, copiada desde
% figures/eda/entradas/01_colinealidad.pdf.
% ⚠️ Esta es la de las 54 columnas. La version REDUCIDA, de 34, va en §\ref{sec:sol_ridge}:
%    las dos juntas enseñan lo que el descarte se lleva por delante.

\subsection{Escalas, asimetrías y ángulos}
\label{sec:escalas}

\begin{itemize}
    \item \textbf{Las escalas son muy diferentes}: $10{,}5$ órdenes de magnitud entre la
          columna más pequeña y la más grande de la tabla, y $8{,}3$ en el vector de nodo.
          \textbf{Decisión: normalizar.}
    \item \textbf{Asimetría fuerte en $34$ columnas}, casi todas recuentos con cola larga a la
          derecha. \textbf{Decisión: logaritmo antes del escalado.}
    \item \textbf{Cinco columnas tienen recorrido intercuartílico nulo}, entre ellas las dos de
          puertas rotas, que valen cero en el $87$\,\% de los circuitos. 
          
          \textbf{Decisión: esas
          cinco no pueden llevar escalado robusto}, porque dividiría por cero.
\end{itemize}

% ⚠️ FORMATO: titulo arriba y fuente abajo (Instrucciones §1.3, p. 6).
% Figura propia, de notebooks/01_eda_gnn/01j, copiada desde
% figures/eda_gnn/nodos/02_escalas.pdf.


% =====================================================================================
\section{Preprocesado}
\label{sec:prep}
% ⚠️ REESCRITO EN CLAVE DE DECISIONES (sep-2026). El razonamiento largo de cada una
%    esta en notebooks/02_preprocesado/DECISIONES_TOMADAS.md.

La regla que se siguió durante todo el preprocesado ha sido:

\begin{center}
    \doublebox{
    \begin{minipage}{0.85\linewidth} \vspace{0.5em}
        \centering
\textit{Toda constante estimada a partir de los datos (medias, escalas, umbrales o cualquier intercepto) se ajusta únicamente con \textbf{\texttt{train}} para evitar la fuga de información (data leakage)}.
        \vspace{0.5em}
    \end{minipage}}
\end{center}

% 🔴 LO QUE ERA EXCLUSIVO DE CADA MODELO SE MOVIO AL CAPITULO 5 (sep-2026, indicacion
%    del tutor): la rama tabular a §5.1 y la del grafo a §5.3, cada una con su modelo.
%    Aqui se queda SOLO lo comun. La figura de colinealidad reducida y la tabla de poda
%    se fueron con ellas. NO reponerlas: estarian en dos sitios.
De esa regla salen dos ramas de preprocesado \textbf{independientes}: una para la matriz
tabular que consumen Ridge y Random Forest, y otra para los grafos que consume la red. Ambas comparten la partición y el umbral, pero no comparten el tratamiento. Cada una se describe junto al modelo que la usa, en los apartados~\ref{sec:sol_ridge} y~\ref{sec:gnn}.

% =====================================================================================
% =====================================================================================

\section{Estrategia de entrenamiento}
\label{sec:entrenamiento}
% =====================================================================================

La estrategia del entrenamiento se compone de los seis peldaños de la
Figura~\ref{fig:peldanos}. Cada uno añade una capacidad sobre el anterior, y la pregunta que
lleva al lado es justamente lo que ese peldaño aísla: si al subirlo la métrica no se mueve,
esa capacidad no aporta nada en este problema.

% ⚠️ FORMATO: titulo arriba y fuente abajo (Instrucciones §1.3, p. 6).
% Figura propia en TikZ. Fuente: figuras/cap06/peldanos_tikz.tex, que lleva dentro las
% medidas y el aviso de que la segunda linea de cada bloque NO puede partirse.
% ⚠️ Esta figura SUSTITUYO a la lista enumerada que habia aqui (sep-2026). No reponer la
%    lista: dirian lo mismo dos veces.
\begin{figure}[htbp]
    \centering
    \caption{Los seis peldaños de la estrategia de entrenamiento.}
    \label{fig:peldanos}

    % =====================================================================================
% La escalera de la estrategia de entrenamiento — figura propia en TikZ.
%
% IDEA: cada peldaño es un BLOQUE que se apoya en el anterior y sobresale hacia la
% derecha. Los bloques se tocan (no hay hueco entre ellos), que es lo que hace que el
% conjunto se lea como una escalera y no como una lista inclinada. Se sube de izquierda
% a derecha porque cada peldaño AÑADE una capacidad sobre el de abajo, y la pregunta en
% gris es justamente lo que ese peldaño aisla frente al anterior.
%
% CODIGO DE COLOR (informativo, no decorativo):
%   azulunir  -> los tres peldaños de REFERENCIA: no son modelos de la comparativa,
%                estan para fijar el suelo contra el que se mide.
%   acentofig -> los TRES MODELOS que se comparan de verdad (§ del capitulo).
%
% ⚠️ ANCHOS MEDIDOS. Caja de texto de la memoria: 15,02 cm.
%    bloque 6,8 cm de texto + 2x0,176 de margen = 7,15 · desplazamiento 1,25 por peldaño
%    -> borde derecho en 5x1,25 + 7,15 = 13,40 · mas la flecha y su rotulo a la izquierda
%    (~0,95) da 14,35 cm. NO ensanchar el bloque sin reducir el desplazamiento.
%
% ⚠️ LA SEGUNDA LINEA DE CADA BLOQUE NO DEBE PARTIRSE. A \scriptsize en 6,8 cm entran
%    ~49 caracteres; las seis estan por debajo. Si se reescribe alguna y pasa de ahi, el
%    bloque crece, deja de medir 1,15 cm y la escalera se descuadra.
%
% ⚠️ SIN TITULO dentro de la imagen: lo pone el \caption de LaTeX (norma del TFE, §1.3).
%
% ⚠️ \shorthandoff{<>} es obligatorio: `babel` en español hace activos < y >, y eso rompe
%    la sintaxis de puntas de flecha de TikZ.
% =====================================================================================
\begingroup
\shorthandoff{<>}%
\hyphenpenalty=10000\exhyphenpenalty=10000\relax

\begin{tikzpicture}[
    >={Stealth[length=2mm]},
    font=\footnotesize,
    peldano/.style = {draw=#1, line width=0.9pt, fill=#1!8, rounded corners=1pt,
                      text width=6.8cm, minimum height=1.15cm, align=left,
                      inner sep=5pt, anchor=south west},
    nota/.style    = {font=\scriptsize, color=gray},
  ]

  % --- los seis peldaños ---------------------------------------------------------------
  % Referencias: fijan el suelo, no entran en la comparativa.
  \node[peldano=azulunir] at (0.00, 0.00)
    {\textbf{1. No mitigar}\\[1pt]
     {\scriptsize\color{gray} el suelo: si no se bate, el método no sirve}};

  \node[peldano=azulunir] at (1.25, 1.15)
    {\textbf{2. Factor de supervivencia constante}\\[1pt]
     {\scriptsize\color{gray} la media de entrenamiento · ¿aporta saberla?}};

  \node[peldano=azulunir] at (2.50, 2.30)
    {\textbf{3. Ridge con una sola variable}\\[1pt]
     {\scriptsize\color{gray} la duración total · ¿aporta saber cuánto dura?}};

  % Los tres que se comparan.
  \node[peldano=acentofig] at (3.75, 3.45)
    {\textbf{4. Ridge con todas las variables}\\[1pt]
     {\scriptsize\color{gray} ¿aporta solo la linealidad?}};

  \node[peldano=acentofig] at (5.00, 4.60)
    {\textbf{5. Random Forest}\\[1pt]
     {\scriptsize\color{gray} ¿aporta la no linealidad?}};

  \node[peldano=acentofig] at (6.25, 5.75)
    {\textbf{6. Red de grafos}\\[1pt]
     {\scriptsize\color{gray} ¿aporta la estructura del circuito?}};

  % --- el eje que dice hacia donde se sube ----------------------------------------------
  \draw[->, azulunir, line width=1pt] (-0.55, 0) -- (-0.55, 6.90)
        node[midway, rotate=90, above=1pt, font=\scriptsize, text=black]
        {capacidad del modelo};

  % --- leyenda, en el hueco que deja la escalera abajo a la derecha ----------------------
  % Va aqui y no arriba a la izquierda por dos razones: arriba quedaba pegada al \caption,
  % y abajo a la derecha rellena el unico hueco grande que deja la escalera.
  \draw[draw=azulunir, line width=0.9pt, fill=azulunir!8]  (7.60, 0.72) rectangle (7.88, 1.00);
  \node[nota, anchor=west] at (7.98, 0.86) {peldaños de referencia: fijan el suelo};

  \draw[draw=acentofig, line width=0.9pt, fill=acentofig!8] (7.60, 0.22) rectangle (7.88, 0.50);
  \node[nota, anchor=west] at (7.98, 0.36) {los tres modelos que se comparan};

\end{tikzpicture}%
\endgroup


    {\footnotesize Fuente: Claude Code \cite{anthropic_claude}.}
\end{figure}

Los tres modelos que se comparan son los tres últimos peldaños y cada uno añade una capacidad sobre el anterior: Ridge es el
\textbf{baseline} y fija el suelo, Random Forest añade no \textbf{linealidad} y está porque es el mejor modelo tabular de los que compara la referencia \cite{liao2024machine}. La red de grafos cambia la representación de los datos, y la diferencia entre ella y los dos primeros es \textbf{el valor de la
estructura que la agregación descarta}.

Cabe recalcar una asimetría presente. Ridge y la red de
grafos se entrenan con una \textbf{pérdida de Huber}, que penaliza los errores grandes de forma lineal y evita que los factores de supervivencia extremos se sesgen el entrenamietno. Sin embargo, \textbf{Random Forest no la admite} y no se emplea.


\section{Métricas y validación}
\label{sec:criterios}
% =====================================================================================
% ⚠️ SECCION SIN SUBAPARTADOS Y ACOTADA A DOS O TRES PAGINAS (sep-2026, decision del
%    autor). Tampoco lleva tablas. No trocearla otra vez ni reponerlas.
%
% 🔴 TRES ETIQUETAS DE AQUI SE CITAN DESDE FUERA, y ahora las tres resuelven a esta misma
%    seccion. NO renombrarlas ni quitarlas:
%      sec:metrica_decide -> capitulo 8      sec:cota -> capitulo 5 (dos veces)
%      sec:reportar       -> capitulo 7
%
% ⚠️ SOLO TRES METRICAS: R², MAE y mejora relativa. Spearman se retiro en sep-2026 por
%    decision del autor, aunque el codigo la siga calculando.
%
% Fuente: notebooks/05_comparativa/METRICAS_Y_EVALUACION.md · src/evaluacion.py
% =====================================================================================
El criterio general está fijado en los objetivos del capítulo~\ref{cap:objetivos}: \textbf{batir
a no mitigar}\footnote{Es decir, lograr que el error de la estimación cuántica corregida ($\langle O\rangle_{\text{mitigado}}$) sea estrictamente menor que el error de la ejecución ruidosa original ($\langle O\rangle_{\text{ruidoso}}$) obtenida directamente del procesador.}, y hacerlo \textbf{fuera de la distribución de entrenamiento}. 

\label{sec:metrica_decide}%

\begin{center}
    \doublebox{
    \begin{minipage}{0.85\linewidth} \vspace{0.5em}
        \centering
        \textit{La métrica principal es el \textbf{coeficiente de determinación sobre el factor
        de supervivencia}, por observable.}
        \vspace{0.5em}
    \end{minipage}}
\end{center}

\begin{equation}
\label{eq:metrica_r}
R^2 = 1 - \frac{\sum_i (f_i - \hat f_i)^2}{\sum_i (f_i - \bar f)^2}.
\end{equation}

Además, para evaluar la desviación absoluta sobre la propia cantidad aprendida ($f$), se emplea el Error Absoluto Medio (MAE):

\begin{equation}
\label{eq:metrica_mae}
\text{MAE}_f = \frac{1}{n}\sum_i \left| f_i - \hat f_i \right| .
\end{equation}

Finalmente, para medir la fidelidad del valor cuántico reconstruido y evaluar la utilidad práctica frente a la línea base del hardware, se utiliza la \textbf{mejora relativa}:

\begin{equation}
\label{eq:mejora_relativa}
\text{Mejora Relativa} = 1 - \frac{\text{MAE}(\langle O\rangle_{\text{mitigado}})}
                                  {\text{MAE}(\langle O\rangle_{\text{ruidoso}})} .
\end{equation}

Esta métrica es fundamental para evaluar hasta qué punto la mitigación ha sido efectiva. Un valor de $1$ corresponde al escenario ideal, en el que el error se corrige por completo. Los valores comprendidos entre $0$ y $1$ indican la proporción de error del hardware que se ha conseguido eliminar. Cuando la métrica toma el valor $0$, significa que el modelo no aporta ninguna mejora respecto a no aplicar mitigación. Por el contrario, los valores negativos revelan que la corrección resulta contraproducente, ya que introduce más error del que logra reducir.

\label{sec:cota}%
La evaluación aplica además la \textbf{cota inferior del factor} deducida en el
apartado~\ref{sec:cotas}: una predicción por debajo de $|\langle O\rangle_{\text{ruidoso}}|$ produciría un valor corregido
imposible, así que antes de reconstruir se sube hasta la cota y se reporta qué porcentaje de
predicciones la tocó.
\begin{equation}
\label{eq:cota}
\hat f_{\text{final}} = \max\left(\hat f,\; |\langle O\rangle_{\text{ruidoso}}|\right),
\end{equation}

\noindent donde $\langle O\rangle_{\text{ruidoso}}$ es el valor esperado medido con ruido.

\label{sec:reportar}%
Por último, todo lo anterior se reporta \textbf{desglosado por los tres ejes de
generalización} (tipo, tamaño y tiempo) descritos en el apartado~\ref{sec:dataset}. 
